% Chương 5. Phụ trách: Minh (kể cả mục Độ nhạy)
% Nguồn: experiments/README.md và experiments/results/*.csv.

\chapter{Thực nghiệm}
\label{ch:thuc-nghiem}

\section{Thiết lập thực nghiệm}
\label{sec:thiet-lap}

\section{Bài toán A}
\label{sec:tn-a}

\subsection{So sánh CP-SAT và PuLP/CBC}
\label{subsec:tn-a-solver}
% experiments/results/perf.csv

\subsection{So sánh chiến lược một pha và hai pha}
\label{subsec:tn-a-chien-luoc}
% experiments/results/perf.csv

\subsection{Khả năng mở rộng và gap}
\label{subsec:tn-a-mo-rong}

\section{Bài toán B}
\label{sec:tn-b}
% experiments/results/perf_b.csv, lab_reserve_b.csv

\section{Thời gian xếp lại khi có sự cố}
\label{sec:tn-xep-lai}
% experiments/results/reschedule.csv

\section{Phân tích độ nhạy theo trọng số mục tiêu}
\label{sec:do-nhay}
% experiments/results/sensitivity.csv
